\chapter{Real- и аномальный допуск минимальной зеркальной пары}
\label{ch:v7-minimal-mirror-pair-real-anomaly-gate}

\section{Проверяемый кандидат}

После закрытия кругового маршрута с $A_{(2)}$ остаётся строгий
первопорядковый кандидат минимальной сырой размерности
\begin{equation}
 X_L=(\mathbb C,\mathbb C,L),\qquad
 Y_R=(\mathbb H,\mathbb C,R).
 \label{eq:v7-mirror-pair-candidate}
\end{equation}
Вместе со стандартными вершинами он замыкает цикл
\begin{equation}
 Q_L\to u_R\to X_L\to e_R\to L_L\to Y_R\to Q_L.
 \label{eq:v7-mirror-pair-cycle}
\end{equation}
Все шесть рёбер меняют хиральность и сохраняют хотя бы одну бимодульную
координату. Следовательно, граф проходит строгий первый порядок.

\section{Real-замыкание не является аномальным замыканием}

В KO-размерности с противоположной хиральностью Real-образ имеет вид
\begin{equation}
 J:(i,j,\chi)\longmapsto(j,i,-\chi),
 \label{eq:v7-mirror-pair-real-map}
\end{equation}
поэтому
\begin{equation}
 JX_L=(\mathbb C,\mathbb C,R),\qquad
 JY_R=(\mathbb C,\mathbb H,L).
 \label{eq:v7-mirror-pair-real-images}
\end{equation}
Эти вершины завершают частице-античастичный Real-носитель. Однако
$J$-удвоение не добавляет независимый набор физических Weyl-полей после
устранения фермионного удвоения и потому само по себе не сокращает аномалии.

\section{Каноническое зеркальное чтение}

Наименьшая пара повторяет калибровочные типы заряженного лептонного
синглета и слабого лептонного дублета, но с противоположными хиральностями:
\begin{equation}
 X_L\sim(\mathbf1,\mathbf1)_{-1},\qquad
 Y_R\sim(\mathbf1,\mathbf2)_{-1/2}.
 \label{eq:v7-mirror-pair-canonical-charges}
\end{equation}
В соглашениях Тома IV её локальные коэффициенты равны
\begin{align}
 \mathcal A_{221}&=-Y(Y_R)=\frac12,\nonumber\\
 \mathcal A_{\mathrm{grav}^2 1}&=Y(X_L)-2Y(Y_R)=0,\nonumber\\
 \mathcal A_{111}&=Y(X_L)^3-2Y(Y_R)^3=-\frac34.
 \label{eq:v7-mirror-pair-canonical-anomalies}
\end{align}
Кроме того, добавлен ровно один физический слабый дублет, так что
\begin{equation}
 N_{\mathbf2}^{\rm new}=1\pmod 2.
 \label{eq:v7-mirror-pair-witten-parity}
\end{equation}
Следовательно, присутствуют и локальная, и глобальная $SU(2)$-аномалии.
Связь аномального сокращения с multiplicity-матрицей и унимодулярностью
конечной геометрии является стандартной частью классификации
\cite{v7mirror-krajewski,v7mirror-alvarez}; нечётное число Weyl-дублетов
даёт глобальное препятствие Виттена \cite{v7mirror-witten}.

\section{Заряд нельзя перенастроить внутри двух вершин}

Пусть $x=Y(X_L)$ и $y=Y(Y_R)$ произвольны. Тогда локальное сокращение требует
\begin{equation}
 -y=0,\qquad x-2y=0,\qquad x^3-2y^3=0,
 \label{eq:v7-mirror-pair-general-anomaly-system}
\end{equation}
откуда $x=y=0$. Но один слабый дублет остаётся нечётным и при нулевом
гиперзаряде. Значит, никакой выбор двух $U(1)$-зарядов не исправляет эту
пару.

\section{Массовая ловушка}

Те же координаты разрешают прямые противоположно-хиральные блоки
\begin{equation}
 X_L\longleftrightarrow e_R,\qquad
 L_L\longleftrightarrow Y_R.
 \label{eq:v7-mirror-pair-direct-masses}
\end{equation}
При общей невырожденной матрице они превращают исходные лёгкие лептоны в
тяжёлые вектороподобные пары. Сохранение наблюдаемого хирального $H_{15}$
потребовало бы вручную задать ядро новой массовой матрицы.

\section{Минимальный консервативный ремонт}

Независимый аномально безопасный сектор требует добавить противоположные
хиральности обоих новых представлений:
\begin{equation}
 X_R\sim(\mathbf1,\mathbf1)_{-1},\qquad
 Y_L\sim(\mathbf1,\mathbf2)_{-1/2}.
 \label{eq:v7-mirror-pair-vectorlike-repair}
\end{equation}
Тогда новый сектор попарно вектороподобен,
\begin{equation}
 \mathcal A_{221}=\mathcal A_{\mathrm{grav}^2 1}
 =\mathcal A_{111}=0,
 \qquad N_{\mathbf2}^{\rm new}=0\pmod2.
 \label{eq:v7-mirror-pair-repair-anomalies}
\end{equation}
Но цена возрастает с двух до четырёх новых физических хиральных вершин, а
их массы и связи ещё не выведены из одного родителя.

\begin{theorem}[Провал минимальной двухвершинной пары]
Пара $X_L\oplus Y_R$ замыкает строгий первопорядковый цикл и допускает
формальное Real-удвоение. Однако она содержит нечётное число слабых дублетов,
не сокращает локальные аномалии и в каноническом чтении допускает прямое
утяжеление исходных лептонов. Поэтому она не является физически допустимым
расширением $H_{15}$.
\end{theorem}

\begin{nogocriterion}[Real не заменяет anomaly]
Нельзя считать пару допустимой только потому, что к ней добавлены
$J$-сопряжённые вершины. Требуется отдельное сокращение аномалий на физическом
Weyl-секторе и сохранение лёгкого хирального спектра без ручного рангового
условия.
\end{nogocriterion}

\begin{protocol}
\begin{enumerate}
 \item строгий первый порядок шестикромочного цикла: \textbf{пройден};
 \item формальное Real-замыкание: \textbf{пройдено};
 \item локальные аномалии канонической пары: \textbf{не сокращаются};
 \item глобальная $SU(2)$-аномалия: \textbf{присутствует};
 \item перенастройка двух гиперзарядов исправляет результат: \textbf{нет};
 \item лёгкий $H_{15}$ сохраняется при общей массе: \textbf{нет};
 \item минимальный аномально безопасный ремонт: \textbf{четыре вершины};
 \item следующий гейт: \textbf{четырёхвершинное вектороподобное замыкание
 против полного зеркального поколения}.
\end{enumerate}
\end{protocol}

Машинный сертификат сохранён в
\path{s2t_v7_minimal_mirror_pair_real_anomaly_gate_results.json}.

\begin{thebibliography}{9}
\bibitem{v7mirror-krajewski}
T.~Krajewski,
\emph{Classification of finite spectral triples},
Journal of Geometry and Physics 28 (1998), 1--30; arXiv:hep-th/9701081.
\bibitem{v7mirror-alvarez}
E.~Alvarez, J.~M.~Gracia-Bondia, C.~P.~Martin,
\emph{Anomaly cancellation and the gauge group of the Standard Model in NCG},
Physics Letters B 364 (1995), 33--40; arXiv:hep-th/9506115.
\bibitem{v7mirror-witten}
E.~Witten,
\emph{An SU(2) anomaly},
Physics Letters B 117 (1982), 324--328.
\end{thebibliography}